% \newcommand*{\StylePath}{../styles}
% % 2014.08.05: This preamble block is now moved to each individual main file, because \import is now used instead of \input
% \documentclass[12pt]{article}
% \usepackage{amsmath,amssymb}
% \usepackage{color,graphicx,fancybox,pause}
% \usepackage[utf8]{vntex} % type Vietnamese
% \usepackage{sagetex} % run Sagemath within Latex

\renewcommand{\familydefault}{cmss}
\renewcommand{\floatpagefraction}{1}

%------------------------------------------
%  PDF
%------------------------------------------
\pdfoutput=1
\pdfcompresslevel=0
\pdfimageresolution=600
\pdfpkresolution=600
\pagestyle{empty}
\pdfpageheight=11.1cm
\pdfpagewidth=14.8cm
%------------------------------------------
% Text dimensions
%------------------------------------------
\textheight=12cm
\textwidth=12.2cm
\topmargin-3cm
\oddsidemargin-1.2cm
\parindent0cm
\parskip0ex
%\arraycolsep0.4ex
%------------------------------------------
%  Lengths
%------------------------------------------
\renewcommand{\baselinestretch}{1.2} % These two lines, although not called in .tex file, affect the output
\renewcommand{\arraystretch}{1}
\newcommand{\spa}[1]{\par\vspace*{#1ex}}
\newcommand{\skipthis}[1]{}

%Set standard display distances
%\rennewcommand{\arraystretch}{1.1}

\newcommand{\sdd}{
\setlength{\jot}{1ex}
\setlength{\abovedisplayskip}{1ex}
\setlength{\abovedisplayshortskip}{1ex}
\setlength{\belowdisplayskip}{1ex}
\setlength{\belowdisplayshortskip}{1ex}}

\newlength{\tew}\setlength{\tew}{\textwidth}
\newlength{\pah}\setlength{\pah}{\pdfpageheight}
\newlength{\storea}
\newlength{\storeb}
\newlength{\storec}
\newlength{\stored}
\newlength{\storee}
\newlength{\storef}

%Scale display distance box
\newcommand{\sdb}[2]{
\storea\abovedisplayskip
\storeb\abovedisplayshortskip
\storec\belowdisplayskip
\stored\belowdisplayshortskip
\setlength{\abovedisplayskip}{#1\abovedisplayskip}
\setlength{\abovedisplayshortskip}{#1\abovedisplayshortskip}
\setlength{\belowdisplayskip}{#1\belowdisplayskip}
\setlength{\belowdisplayshortskip}{#1\belowdisplayshortskip}

#2

\setlength{\abovedisplayskip}{\storea}
\setlength{\abovedisplayshortskip}{\storeb}
\setlength{\belowdisplayskip}{\storec}
\setlength{\belowdisplayshortskip}{\stored}
}

%Scale display distance box with different factors top and bottom
\newcommand{\sdbe}[3]{
\storea=\abovedisplayskip
\storeb=\abovedisplayshortskip
\storec=\belowdisplayskip
\stored=\belowdisplayshortskip
\setlength{\abovedisplayskip}{#1\abovedisplayskip}
\setlength{\abovedisplayshortskip}{#1\abovedisplayshortskip}
\setlength{\belowdisplayskip}{#2\belowdisplayskip}
\setlength{\belowdisplayshortskip}{#2\belowdisplayshortskip}

#3

\setlength{\abovedisplayskip}{\storea}
\setlength{\abovedisplayshortskip}{\storeb}
\setlength{\belowdisplayskip}{\storec}
\setlength{\belowdisplayshortskip}{\stored}
}

%Baselineskip box: scale baslineskip
\newcommand{\blb}[2]{
\storee=\baselineskip
\setlength{\baselineskip}{#1\baselineskip}
#2

\setlength{\baselineskip}{\storee}
}

%------------------------------------------
\newcommand{\enu}[1]{ \begin{enumerate} #1 \end{enumerate} }
\newcommand{\ite}[1]{\vspace*{0ex}\begin{list}{$\bullet$}{\leftmargin2.5ex}#1\end{list} } % itemize list without indentation
\newcommand{\itemi}[1]{\begin{itemize}#1\end{itemize} }

\renewcommand{\labelenumi}{\arabic{enumi}.}
\renewcommand{\labelenumii}{(\alph{enumii})}

\newcommand{\aln}[1]{\noindent\begin{align*}#1\end{align*}}
\newcommand{\mun}[1]{\noindent\begin{multline*}#1\end{multline*}}
\newcommand{\eqn}[1]{\spa{-3}$$#1$$\spa{0}}
\newcommand{\eqns}[1]{\spa{-2}$$#1$$\spa{0}}
\newcommand{\gan}[1]{\begin{gather*}#1\end{gather*}}
\newcommand{\ear}[1]{\begin{eqnarray}#1\end{eqnarray}}
\newcommand{\ean}[1]{\begin{eqnarray*}#1\end{eqnarray*}}

\newcommand{\arr}[2]{\begin{array}{#1}#2\end{array}}
\newcommand{\tab}[2]{\begin{tabular}{#1}#2\end{tabular}}
\newcommand{\mas}[2]{\left[\begin{array}{#1}#2\end{array}\right]}
\newcommand{\mat}[2]{\left(\begin{array}{#1}#2\end{array}\right)}
%\newcommand{\mat}[2]{\mas{#1}{#2}}
\newcommand{\matl}[4]{\left#1\begin{array}{#2}#3\end{array}\right#4}
\newcommand{\dmas}[2]{
\mas{ccc}{#1&&0\\[-1.5ex]&\ddots&\\[-1.5ex]0&&#2}}
\newcommand{\quo}[1]{\noindent\vspace*{-1ex}\begin{quote}#1\end{quote}}
\newcommand{\cen}[1]{\begin{center}#1\end{center}}
\newcommand{\epro}{ \mbox{\rule{2mm}{2mm}} }

%-----------------------------------------------------------------------
%title frontpage CTUT and STW
\newcommand{\titfstw}[2]{%\vspace*{-1ex}
{\hspace*{-0.0\tew}{\bf\Large\color{black}{{#1}}}\\[-1.6ex] }\\[0.70\pah]
\fboxsep=0.3ex
\hspace*{-0.12\tew}
\colorbox{mag}{
\begin{minipage}{1.25\tew}
\tiny  \hfill \phantom{\arabic{page}/\pageref{lastpage}}\hspace*{1cm}
\end{minipage}}
\hspace*{-0.12\tew}
\begin{minipage}{1.25\tew}
\vspace*{-0.7ex}
\mbox{}\hspace{0.11\tew}\footnotesize\phantom{\includegraphics[width=1.5cm]{logo_ctut}\hfill\raisebox{1ex}{Doãn Minh Đăng}}\hfill
\includegraphics[width=1.5cm]{logo_ctut}\hspace*{0.8cm}
\end{minipage}\\[-0.88\pah]

{\large #2}}
%-----------------------------------------------------------------------
%title CTUT and STW
\newcommand{\titstw}[1]{%\vspace*{-1ex}
{\hspace*{-0.0\tew}{\bf\large\color{black}{{#1}}}\\[-1.6ex] }\\[0.70\pah]
\fboxsep=0.3ex
\hspace*{-0.12\tew}
\colorbox{mag}{
\begin{minipage}{1.25\tew}
\tiny  \hfill \arabic{page}/\pageref{lastpage}\hspace*{1cm}
\end{minipage}}
\hspace*{-0.12\tew}
\begin{minipage}{1.25\tew}
\vspace*{-0.7ex}
\mbox{}\hspace{0.02\tew}
\raisebox{0.5ex}{\includegraphics[width=1.2cm]{logo_ctut}}
\hfill\footnotesize\raisebox{1.2ex}{Doãn Minh Đăng}\hfill
\raisebox{0.1ex}{\includegraphics[width=1.5cm]{logo_ctut}}\hspace*{0.8cm}
\end{minipage}\\[-0.88\pah]}
%-----------------------------------------------------------------------
%title CTUT logo small
\newcommand{\titls}[1]{%\vspace*{-1ex}
{\hspace*{-0.0\tew}{\bf\large\color{black}{{#1}}}\\[-1.6ex] }\\[0.70\pah]
\fboxsep=0.3ex
\hspace*{-0.12\tew}
\colorbox{mag}{
\begin{minipage}{1.25\tew}
\tiny  \hfill \arabic{page}/\pageref{lastpage}\hspace*{1cm}
\end{minipage}}
\hspace*{-0.12\tew}
\begin{minipage}{1.25\tew}
\vspace*{-0.7ex}
% \mbox{}\hspace{0.11\tew}\footnotesize\raisebox{1ex}{Carsten Scherer\ \ }\includegraphics[width=0.6cm]{logo_ctut}
% \hspace{26ex}\raisebox{1ex}{Siep Weiland\ \ }\includegraphics[width=0.6cm]{logo_ctut}\hspace*{0.8cm}
\mbox{}\hspace{0.11\tew}\footnotesize\raisebox{1ex}{Doãn Minh Đăng}\hfill
\includegraphics[width=0.5cm]{logo_ctut}\hspace*{0.8cm}
\end{minipage}\\[-0.88\pah]}
%-----------------------------------------------------------------------
%title underlined left
\newcommand{\titline}[1]{{
\hspace*{-0.0\tew}\hfill{\bf \color[rgb]{0,0,0}{#1}}\\[-1.6ex]}
\hspace*{-0.0\tew}\color[rgb]{0,0,0}{\rule{1.00\tew}{.12ex}}\\[-2ex]
\marginpar{{\tiny \vspace*{9cm}\hspace*{0.55cm}\arabic{page}/\pageref{lastpage}}}}

%-----------------------------------------------------------------------
%just title
\newcommand{\tit}[1]{
{\hspace*{-0.0\tew}{\bf\large \color{black}{{#1}}}\\[-1.6ex] }\\[0.79\pah]
\begin{minipage}{1.09\tew}
\tiny  \hfill \arabic{page}/\pageref{lastpage}
\end{minipage}\\[-0.9\pah]
}

%-----------------------------------------------------------------------
%title CTUT logo large
\newcommand{\titlogo}[2]{%\vspace*{-1ex}
{\hspace*{-0.0\tew}{\bf\Large\color{black}{{#1}}}\\[-1.6ex] }\\[0.70\pah]
\fboxsep=0.3ex
\hspace*{-0.12\tew}
\colorbox{mag}{
\begin{minipage}{1.25\tew}
\tiny  \hfill \arabic{page}/\pageref{lastpage}\hspace*{1cm}
\end{minipage}}
\hspace*{-0.12\tew}
\begin{minipage}{1.25\tew}
\vspace*{-0.7ex}
\mbox{}\hspace{0.11\tew}\footnotesize\raisebox{1ex}{Doãn Minh Đăng}\hfill
\includegraphics[width=0.6cm]{logo_ctut}\hspace*{0.8cm}
\end{minipage}\\[-0.88\pah]

{\large #2}}
%-----------------------------------------------------------------------
%title first page CTUT logo large
\newcommand{\titf}[2]{%\vspace*{-1ex}
{\hspace*{-0.0\tew}{\bf\Large\color{blue}{{#1}}}\\[-1.6ex] }\\[0.70\pah]
\fboxsep=0.3ex
\hspace*{-0.12\tew}
\colorbox{mag}{
\begin{minipage}{1.25\tew}
\tiny  \hfill \phantom{\arabic{page}/\pageref{lastpage}}\hspace*{1cm}
\end{minipage}}
\hspace*{-0.12\tew}
\begin{minipage}{1.25\tew}
\vspace*{-0.7ex}
\mbox{}\hspace{0.11\tew}\footnotesize\raisebox{1ex}{Doãn Minh Đăng}\hfill
\includegraphics[width=0.5cm]{logo_ctut}\hspace*{0.8cm}
\end{minipage}\\[-0.88\pah]

{\large #2}}

%-----------------------------------------------------------------------
% Boxes and minipages
%-----------------------------------------------------------------------
\newcommand{\sbo}[1]{
\shadowbox{
\vspace{1ex}
\begin{minipage}{0.97\tew}
\vspace*{1.5ex}
#1
\par\vspace{1ex}
\end{minipage}}}


\newcommand{\cpa}[3]{
{\fboxsep=1.2ex
\colorbox{#1}{
\begin{minipage}{#2}
#3
\end{minipage}}}
}

\newcommand{\cbo}[2]{
\fboxsep=1ex
\colorbox{#1}{
#2}}

\newcommand{\cbow}[3]{
\fboxsep=#2ex
\colorbox{#1}{
#3}}

\newcommand{\mini}[2]{ {\begin{minipage}[t]{#1}\sloppy #2 \end{minipage} } }

\newcommand{\twomini}[6]{
\setlength{\storee}{0.975\tew}
\addtolength{\storee}{-#1}\addtolength{\storee}{-#3}
\addtolength{\storee}{#2}
\mini{#1}{#5}\hspace{\storee}\raisebox{#4}{\mini{#3}{#6}} }
%\hspace{\storee\raisebox{#2}

\newcommand{\onemini}[3]{\hspace{#1}\mini{#2}{#3}}

%%%%%%%%%%%%%%%%%%%%%%%%%%%%%%%%%%%%%%%%%%%%%%%%%%%%%%%%%%%%%%%%%%%%%%%%%%%%%%%%%%%%%%%%%%%%%%%%%%%%%%%%%%%%%%%%%%%
%Colors
\newcommand{\red}[1]{{\color[rgb]{1,0,0}#1}}
\newcommand{\green}[1]{{\color[rgb]{0,0.7,0}#1}}
\newcommand{\blue}[1]{{\color[rgb]{0,0,1}#1}}
\newcommand{\dblue}[1]{{\color[rgb]{0,0,.4}#1}}
%\newcommand{\blue}[1]{{\color[rgb]{0,0.5,.8}#1}}
\newcommand{\black}[1]{{\color[rgb]{0,0,0}#1}}
%\newcommand{\yellow}[1]{{\color[rgb]{1,1,0}#1}}
\newcommand{\yellow}[1]{{\color[rgb]{1,0.95,.5}#1}}
%TUD RBG   0   174   239
\definecolor{lmag}{rgb}{0,0.82,1}
% \definecolor{mag}{rgb}{0,0.70,.88}
\newcommand{\tmag}[1]{{\color[rgb]{0,0.60,.78}#1}}
% \definecolor{mag}{rgb}{0,0.75,.93}
% \definecolor{mag}{rgb}{0,0.6,.8} %original
\definecolor{mag}{rgb}{0,0.7,.9}
\definecolor{lblue}{rgb}{0.6,0.8,1}
\definecolor{lgreen}{rgb}{0,0.92,0.7}
\definecolor{lred}{rgb}{1,0.75,0.75}
%\definecolor{lred}{rgb}{1,0.80,0.80}
\definecolor{lblue}{rgb}{0.6,0.8,1}
\definecolor{lgrey}{rgb}{0.85,0.85,0.85}

\definecolor{lye}{rgb}{1,0.9,0.6}
%\definecolor{lye}{rgb}{1,0.95,0.7}
%\definecolor{lgr}{rgb}{1,0.9,0.6}

\newcommand{\ypa}[1]{\cpa{lye}{0.95\tew}{#1}}
\newcommand{\bpa}[1]{\cpa{lblue}{0.95\tew}{#1}}
\newcommand{\rpa}[1]{\cpa{lred}{0.95\tew}{#1}}
%%%%%%%%%%%%%%%%%%%%%%%%%%%%%%%%%%%%%%%%%%%%%%%%%%%%%%%%%%%%%%%%%%%%%%%%%%%%%%%%%%%%%%%%%%%%%%%%%%%%%%%%%%%%%%%%%%%

\newcommand{\minimize}[2]{
\begin{array}{ll}
\mbox{minimize}    &\ \  #1 \\
\mbox{subject\ to} &\ \  #2
\end{array} }
\newcommand{\opt}{\text{opt}}

\newcommand{\X}{\red{X}}
\newcommand{\pa}{\blue{\delta}}
\newcommand{\ve}{\blue{v}}
\newcommand{\pas}{\blue{\bm{\delta}}}
\newcommand{\ves}{\bm{\blue{v}}}
\newcommand{\pad}{\blue{\bm{\dot{\delta}}}}
\newcommand{\pam}{\blue{\Delta(\pa)}}

\newcommand{\yc}{\blue{y}}
\newcommand{\cge}{\succcurlyeq}
\newcommand{\cle}{\preccurlyeq}
\newcommand{\cl}{\prec}
\newcommand{\cg}{\succ}
\newcommand{\rk}{\text{rank}}
\newcommand{\col}{\text{col}}

%\newcommand{\pa}{\blue{\delta}}
%\newcommand{\setpa}{\blue{\bm{\delta}}}

\renewcommand{\c}[1]{ {\cal #1} }
\newcommand{\bm}[1]{\text{\boldmath$ #1 $} }
\newcommand{\dis}{\displaystyle}
\newcommand{\ab}[1]{\langle #1 \rangle}
\newcommand{\ub}[2]{\underbrace{#1}_{#2}}

\newcommand{\diag}{ \operatornamewithlimits{diag} }
\newcommand{\re}{ \operatorname{Re} }
\newcommand{\he}{ \operatorname{He} }
\renewcommand{\t}[1]{ \tilde{\! #1\,} }
\newcommand{\tc}[1]{ \tilde{{\cal #1}} }
% \renewcommand{\h}[1]{ \hat{#1} } % 2014.08.06: remove this short-hand command, because it cause conflict with Vietnamese language
% If \h is used, some characters with accent like ể will cause error for amsmath: ``\hat allowed only in math mode''
\newcommand{\hc}[1]{ \hat{{\cal #1}} }
\newcommand{\esup}{ \operatornamewithlimits{ess\,sup} }
\newcommand{\tr}{ \operatorname{trace} }
\newcommand{\co}{ \mbox{co} }
\renewcommand{\i}{i}
\newcommand{\io}{i\omega}
\newcommand{\om}{\omega}
\newcommand{\la}{\lambda}
\newcommand{\ga}{\gamma}

\newcommand{\Cl}{ {\mathbb{C}^-} }
\newcommand{\Ci}{ {\mathbb{C}^0} }
\newcommand{\Cr}{ {\mathbb{C}^+} }
\newcommand{\C}{ {\mathbb{C}} }
\newcommand{\Cset}{ {\mathbb{C}} }

\newcommand{\Q}{ {\mathbb{Q}} }
\newcommand{\R}{ {\mathbb{R}} }
\newcommand{\Rset}{ {\mathbb{R}} }
\renewcommand{\H}{ {\mathbb{H}} }
\renewcommand{\S}{ {\mathbb{S}} }
\newcommand{\Nset}{ {\mathbb{N}} }
\newcommand{\Zset}{ {\mathbb{Z}} }
\newcommand{\alert}[1]{{\bf #1}}
%%%%%%%%%%%%%%%%%%%%%%%%%%%%%%%%%%%%%%%%%%%%%%%%%%%%%%%%%%%%%%%%%%%%%%%%%%%%%%%%%%%%%%%%%%%%%%%%%%%%%%%%%%%%%%%%%%%

\newcommand{\Alf}{{D}}
\newcommand{\Blf}{{C}}
\newcommand{\Clf}{{B}}
\newcommand{\Dlf}{{A}}


\newcommand{\Ac}{\green{A_c}}
\newcommand{\Bc}{\green{B_c}}
\newcommand{\Cc}{\green{C_c}}
\newcommand{\Dc}{\green{D_c}}

\newcommand{\Acl}{{\cal A}}
\newcommand{\Bcl}{{\cal B}}
\newcommand{\Ccl}{{\cal C}}
\newcommand{\Dcl}{{\cal D}}
\newcommand{\Xcl}{{ \green{\cal X} }}
\newcommand{\Ycl}{{\cal Y}}
\newcommand{\Zcl}{{\cal Z}}
\newcommand{\Pcl}{{\cal P}}
%closed-loop transformed
\newcommand{\A}{\bm{A}}
\newcommand{\B}{\bm{B}}
%\renewcommand{\C}{\bm{C}}
\newcommand{\D}{\bm{D}}
\renewcommand{\X}{\bm{X}}
\newcommand{\DX}{\bm{Y}}

\renewcommand{\v}{\red{v}}
\newcommand{\K}{\red{K}}
\renewcommand{\L}{\red{L}}
\newcommand{\M}{\red{M}}
\newcommand{\N}{\red{N}}
\renewcommand{\X}{\red{X}}
\newcommand{\Y}{\red{Y}}
\newcommand{\Z}{\red{Z}}
\newcommand{\F}{\blue{F}}

\newcommand{\T}{\c{T}}

\newcommand{\Ri}{\R\cup\{\infty\}}

%%%% Dang added %%%%
\newcommand{\xb}{\mathbf{x}}
\newcommand{\grad}{\nabla}
\newtheorem{Lem}{Bổ đề}
\newtheorem{Rem}{Lưu ý}
\newtheorem{Def}{Định nghĩa}
\newtheorem{Thm}{Định lý}
\newtheorem{Exa}{Ví dụ}



\documentclass[12pt]{article}
% \usepackage[utf8]{vietnam} % type Vietnamese
\usepackage[utf8]{inputenc}
\usepackage[vietnam]{babel}
\usepackage{amsmath,amsfonts,amssymb}
\usepackage{color,graphicx,fancybox} % pause.sty cung cap chuc nang pause giong nhu o documentclass beamer, tuy nhien dung o documentclass article nay thi khong co tac dung, nen bo di
\usepackage{sagetex} % run Sagemath within Latex
\usepackage{import} % 2014.08.05 so that using \import command instead of \input in the main files
% because \input or \include only search for the files with relative path to the MAIN tex files,
% while \import allows the figures inside the ``styles'' folder to be searched as well.
% see: http://sumitrangwala.wordpress.com/2009/03/09/latex-tips-and-tricks/

\graphicspath{{../styles/}} % khong co graphicspath thi file ../styles/slides.tex khong search duoc cac file logo, although da dung subimport
% Dung graphicspath la giai phap temporary, vi neu trong slides.tex ma co \input{otherfile.tex} thi no khong search duoc qua relative path
% Khi do phai kiem tra lam sao su dung duoc \subimport.

% \subimport*{../styles/}{slides} % dung thu ngay 2014.08.05 khong thay tac dung.

% 2014.08.05: This preamble block is now moved to each individual main file, because \import is now used instead of \input
% \documentclass[12pt]{article}
% \usepackage{amsmath,amssymb}
% \usepackage{color,graphicx,fancybox,pause}
% \usepackage[utf8]{vntex} % type Vietnamese
% \usepackage{sagetex} % run Sagemath within Latex

\renewcommand{\familydefault}{cmss}
\renewcommand{\floatpagefraction}{1}

%------------------------------------------
%  PDF
%------------------------------------------
\pdfoutput=1
\pdfcompresslevel=0
\pdfimageresolution=600
\pdfpkresolution=600
\pagestyle{empty}
\pdfpageheight=11.1cm
\pdfpagewidth=14.8cm
%------------------------------------------
% Text dimensions
%------------------------------------------
\textheight=12cm
\textwidth=12.2cm
\topmargin-3cm
\oddsidemargin-1.2cm
\parindent0cm
\parskip0ex
%\arraycolsep0.4ex
%------------------------------------------
%  Lengths
%------------------------------------------
\renewcommand{\baselinestretch}{1.2} % These two lines, although not called in .tex file, affect the output
\renewcommand{\arraystretch}{1}
\newcommand{\spa}[1]{\par\vspace*{#1ex}}
\newcommand{\skipthis}[1]{}

%Set standard display distances
%\rennewcommand{\arraystretch}{1.1}

\newcommand{\sdd}{
\setlength{\jot}{1ex}
\setlength{\abovedisplayskip}{1ex}
\setlength{\abovedisplayshortskip}{1ex}
\setlength{\belowdisplayskip}{1ex}
\setlength{\belowdisplayshortskip}{1ex}}

\newlength{\tew}\setlength{\tew}{\textwidth}
\newlength{\pah}\setlength{\pah}{\pdfpageheight}
\newlength{\storea}
\newlength{\storeb}
\newlength{\storec}
\newlength{\stored}
\newlength{\storee}
\newlength{\storef}

%Scale display distance box
\newcommand{\sdb}[2]{
\storea\abovedisplayskip
\storeb\abovedisplayshortskip
\storec\belowdisplayskip
\stored\belowdisplayshortskip
\setlength{\abovedisplayskip}{#1\abovedisplayskip}
\setlength{\abovedisplayshortskip}{#1\abovedisplayshortskip}
\setlength{\belowdisplayskip}{#1\belowdisplayskip}
\setlength{\belowdisplayshortskip}{#1\belowdisplayshortskip}

#2

\setlength{\abovedisplayskip}{\storea}
\setlength{\abovedisplayshortskip}{\storeb}
\setlength{\belowdisplayskip}{\storec}
\setlength{\belowdisplayshortskip}{\stored}
}

%Scale display distance box with different factors top and bottom
\newcommand{\sdbe}[3]{
\storea=\abovedisplayskip
\storeb=\abovedisplayshortskip
\storec=\belowdisplayskip
\stored=\belowdisplayshortskip
\setlength{\abovedisplayskip}{#1\abovedisplayskip}
\setlength{\abovedisplayshortskip}{#1\abovedisplayshortskip}
\setlength{\belowdisplayskip}{#2\belowdisplayskip}
\setlength{\belowdisplayshortskip}{#2\belowdisplayshortskip}

#3

\setlength{\abovedisplayskip}{\storea}
\setlength{\abovedisplayshortskip}{\storeb}
\setlength{\belowdisplayskip}{\storec}
\setlength{\belowdisplayshortskip}{\stored}
}

%Baselineskip box: scale baslineskip
\newcommand{\blb}[2]{
\storee=\baselineskip
\setlength{\baselineskip}{#1\baselineskip}
#2

\setlength{\baselineskip}{\storee}
}

%------------------------------------------
\newcommand{\enu}[1]{ \begin{enumerate} #1 \end{enumerate} }
\newcommand{\ite}[1]{\vspace*{0ex}\begin{list}{$\bullet$}{\leftmargin2.5ex}#1\end{list} } % itemize list without indentation
\newcommand{\itemi}[1]{\begin{itemize}#1\end{itemize} }

\renewcommand{\labelenumi}{\arabic{enumi}.}
\renewcommand{\labelenumii}{(\alph{enumii})}

\newcommand{\aln}[1]{\noindent\begin{align*}#1\end{align*}}
\newcommand{\mun}[1]{\noindent\begin{multline*}#1\end{multline*}}
\newcommand{\eqn}[1]{\spa{-3}$$#1$$\spa{0}}
\newcommand{\eqns}[1]{\spa{-2}$$#1$$\spa{0}}
\newcommand{\gan}[1]{\begin{gather*}#1\end{gather*}}
\newcommand{\ear}[1]{\begin{eqnarray}#1\end{eqnarray}}
\newcommand{\ean}[1]{\begin{eqnarray*}#1\end{eqnarray*}}

\newcommand{\arr}[2]{\begin{array}{#1}#2\end{array}}
\newcommand{\tab}[2]{\begin{tabular}{#1}#2\end{tabular}}
\newcommand{\mas}[2]{\left[\begin{array}{#1}#2\end{array}\right]}
\newcommand{\mat}[2]{\left(\begin{array}{#1}#2\end{array}\right)}
%\newcommand{\mat}[2]{\mas{#1}{#2}}
\newcommand{\matl}[4]{\left#1\begin{array}{#2}#3\end{array}\right#4}
\newcommand{\dmas}[2]{
\mas{ccc}{#1&&0\\[-1.5ex]&\ddots&\\[-1.5ex]0&&#2}}
\newcommand{\quo}[1]{\noindent\vspace*{-1ex}\begin{quote}#1\end{quote}}
\newcommand{\cen}[1]{\begin{center}#1\end{center}}
\newcommand{\epro}{ \mbox{\rule{2mm}{2mm}} }

%-----------------------------------------------------------------------
%title frontpage CTUT and STW
\newcommand{\titfstw}[2]{%\vspace*{-1ex}
{\hspace*{-0.0\tew}{\bf\Large\color{black}{{#1}}}\\[-1.6ex] }\\[0.70\pah]
\fboxsep=0.3ex
\hspace*{-0.12\tew}
\colorbox{mag}{
\begin{minipage}{1.25\tew}
\tiny  \hfill \phantom{\arabic{page}/\pageref{lastpage}}\hspace*{1cm}
\end{minipage}}
\hspace*{-0.12\tew}
\begin{minipage}{1.25\tew}
\vspace*{-0.7ex}
\mbox{}\hspace{0.11\tew}\footnotesize\phantom{\includegraphics[width=1.5cm]{logo_ctut}\hfill\raisebox{1ex}{Doãn Minh Đăng}}\hfill
\includegraphics[width=1.5cm]{logo_ctut}\hspace*{0.8cm}
\end{minipage}\\[-0.88\pah]

{\large #2}}
%-----------------------------------------------------------------------
%title CTUT and STW
\newcommand{\titstw}[1]{%\vspace*{-1ex}
{\hspace*{-0.0\tew}{\bf\large\color{black}{{#1}}}\\[-1.6ex] }\\[0.70\pah]
\fboxsep=0.3ex
\hspace*{-0.12\tew}
\colorbox{mag}{
\begin{minipage}{1.25\tew}
\tiny  \hfill \arabic{page}/\pageref{lastpage}\hspace*{1cm}
\end{minipage}}
\hspace*{-0.12\tew}
\begin{minipage}{1.25\tew}
\vspace*{-0.7ex}
\mbox{}\hspace{0.02\tew}
\raisebox{0.5ex}{\includegraphics[width=1.2cm]{logo_ctut}}
\hfill\footnotesize\raisebox{1.2ex}{Doãn Minh Đăng}\hfill
\raisebox{0.1ex}{\includegraphics[width=1.5cm]{logo_ctut}}\hspace*{0.8cm}
\end{minipage}\\[-0.88\pah]}
%-----------------------------------------------------------------------
%title CTUT logo small
\newcommand{\titls}[1]{%\vspace*{-1ex}
{\hspace*{-0.0\tew}{\bf\large\color{black}{{#1}}}\\[-1.6ex] }\\[0.70\pah]
\fboxsep=0.3ex
\hspace*{-0.12\tew}
\colorbox{mag}{
\begin{minipage}{1.25\tew}
\tiny  \hfill \arabic{page}/\pageref{lastpage}\hspace*{1cm}
\end{minipage}}
\hspace*{-0.12\tew}
\begin{minipage}{1.25\tew}
\vspace*{-0.7ex}
% \mbox{}\hspace{0.11\tew}\footnotesize\raisebox{1ex}{Carsten Scherer\ \ }\includegraphics[width=0.6cm]{logo_ctut}
% \hspace{26ex}\raisebox{1ex}{Siep Weiland\ \ }\includegraphics[width=0.6cm]{logo_ctut}\hspace*{0.8cm}
\mbox{}\hspace{0.11\tew}\footnotesize\raisebox{1ex}{Doãn Minh Đăng}\hfill
\includegraphics[width=0.5cm]{logo_ctut}\hspace*{0.8cm}
\end{minipage}\\[-0.88\pah]}
%-----------------------------------------------------------------------
%title underlined left
\newcommand{\titline}[1]{{
\hspace*{-0.0\tew}\hfill{\bf \color[rgb]{0,0,0}{#1}}\\[-1.6ex]}
\hspace*{-0.0\tew}\color[rgb]{0,0,0}{\rule{1.00\tew}{.12ex}}\\[-2ex]
\marginpar{{\tiny \vspace*{9cm}\hspace*{0.55cm}\arabic{page}/\pageref{lastpage}}}}

%-----------------------------------------------------------------------
%just title
\newcommand{\tit}[1]{
{\hspace*{-0.0\tew}{\bf\large \color{black}{{#1}}}\\[-1.6ex] }\\[0.79\pah]
\begin{minipage}{1.09\tew}
\tiny  \hfill \arabic{page}/\pageref{lastpage}
\end{minipage}\\[-0.9\pah]
}

%-----------------------------------------------------------------------
%title CTUT logo large
\newcommand{\titlogo}[2]{%\vspace*{-1ex}
{\hspace*{-0.0\tew}{\bf\Large\color{black}{{#1}}}\\[-1.6ex] }\\[0.70\pah]
\fboxsep=0.3ex
\hspace*{-0.12\tew}
\colorbox{mag}{
\begin{minipage}{1.25\tew}
\tiny  \hfill \arabic{page}/\pageref{lastpage}\hspace*{1cm}
\end{minipage}}
\hspace*{-0.12\tew}
\begin{minipage}{1.25\tew}
\vspace*{-0.7ex}
\mbox{}\hspace{0.11\tew}\footnotesize\raisebox{1ex}{Doãn Minh Đăng}\hfill
\includegraphics[width=0.6cm]{logo_ctut}\hspace*{0.8cm}
\end{minipage}\\[-0.88\pah]

{\large #2}}
%-----------------------------------------------------------------------
%title first page CTUT logo large
\newcommand{\titf}[2]{%\vspace*{-1ex}
{\hspace*{-0.0\tew}{\bf\Large\color{blue}{{#1}}}\\[-1.6ex] }\\[0.70\pah]
\fboxsep=0.3ex
\hspace*{-0.12\tew}
\colorbox{mag}{
\begin{minipage}{1.25\tew}
\tiny  \hfill \phantom{\arabic{page}/\pageref{lastpage}}\hspace*{1cm}
\end{minipage}}
\hspace*{-0.12\tew}
\begin{minipage}{1.25\tew}
\vspace*{-0.7ex}
\mbox{}\hspace{0.11\tew}\footnotesize\raisebox{1ex}{Doãn Minh Đăng}\hfill
\includegraphics[width=0.5cm]{logo_ctut}\hspace*{0.8cm}
\end{minipage}\\[-0.88\pah]

{\large #2}}

%-----------------------------------------------------------------------
% Boxes and minipages
%-----------------------------------------------------------------------
\newcommand{\sbo}[1]{
\shadowbox{
\vspace{1ex}
\begin{minipage}{0.97\tew}
\vspace*{1.5ex}
#1
\par\vspace{1ex}
\end{minipage}}}


\newcommand{\cpa}[3]{
{\fboxsep=1.2ex
\colorbox{#1}{
\begin{minipage}{#2}
#3
\end{minipage}}}
}

\newcommand{\cbo}[2]{
\fboxsep=1ex
\colorbox{#1}{
#2}}

\newcommand{\cbow}[3]{
\fboxsep=#2ex
\colorbox{#1}{
#3}}

\newcommand{\mini}[2]{ {\begin{minipage}[t]{#1}\sloppy #2 \end{minipage} } }

\newcommand{\twomini}[6]{
\setlength{\storee}{0.975\tew}
\addtolength{\storee}{-#1}\addtolength{\storee}{-#3}
\addtolength{\storee}{#2}
\mini{#1}{#5}\hspace{\storee}\raisebox{#4}{\mini{#3}{#6}} }
%\hspace{\storee\raisebox{#2}

\newcommand{\onemini}[3]{\hspace{#1}\mini{#2}{#3}}

%%%%%%%%%%%%%%%%%%%%%%%%%%%%%%%%%%%%%%%%%%%%%%%%%%%%%%%%%%%%%%%%%%%%%%%%%%%%%%%%%%%%%%%%%%%%%%%%%%%%%%%%%%%%%%%%%%%
%Colors
\newcommand{\red}[1]{{\color[rgb]{1,0,0}#1}}
\newcommand{\green}[1]{{\color[rgb]{0,0.7,0}#1}}
\newcommand{\blue}[1]{{\color[rgb]{0,0,1}#1}}
\newcommand{\dblue}[1]{{\color[rgb]{0,0,.4}#1}}
%\newcommand{\blue}[1]{{\color[rgb]{0,0.5,.8}#1}}
\newcommand{\black}[1]{{\color[rgb]{0,0,0}#1}}
%\newcommand{\yellow}[1]{{\color[rgb]{1,1,0}#1}}
\newcommand{\yellow}[1]{{\color[rgb]{1,0.95,.5}#1}}
%TUD RBG   0   174   239
\definecolor{lmag}{rgb}{0,0.82,1}
% \definecolor{mag}{rgb}{0,0.70,.88}
\newcommand{\tmag}[1]{{\color[rgb]{0,0.60,.78}#1}}
% \definecolor{mag}{rgb}{0,0.75,.93}
% \definecolor{mag}{rgb}{0,0.6,.8} %original
\definecolor{mag}{rgb}{0,0.7,.9}
\definecolor{lblue}{rgb}{0.6,0.8,1}
\definecolor{lgreen}{rgb}{0,0.92,0.7}
\definecolor{lred}{rgb}{1,0.75,0.75}
%\definecolor{lred}{rgb}{1,0.80,0.80}
\definecolor{lblue}{rgb}{0.6,0.8,1}
\definecolor{lgrey}{rgb}{0.85,0.85,0.85}

\definecolor{lye}{rgb}{1,0.9,0.6}
%\definecolor{lye}{rgb}{1,0.95,0.7}
%\definecolor{lgr}{rgb}{1,0.9,0.6}

\newcommand{\ypa}[1]{\cpa{lye}{0.95\tew}{#1}}
\newcommand{\bpa}[1]{\cpa{lblue}{0.95\tew}{#1}}
\newcommand{\rpa}[1]{\cpa{lred}{0.95\tew}{#1}}
%%%%%%%%%%%%%%%%%%%%%%%%%%%%%%%%%%%%%%%%%%%%%%%%%%%%%%%%%%%%%%%%%%%%%%%%%%%%%%%%%%%%%%%%%%%%%%%%%%%%%%%%%%%%%%%%%%%

\newcommand{\minimize}[2]{
\begin{array}{ll}
\mbox{minimize}    &\ \  #1 \\
\mbox{subject\ to} &\ \  #2
\end{array} }
\newcommand{\opt}{\text{opt}}

\newcommand{\X}{\red{X}}
\newcommand{\pa}{\blue{\delta}}
\newcommand{\ve}{\blue{v}}
\newcommand{\pas}{\blue{\bm{\delta}}}
\newcommand{\ves}{\bm{\blue{v}}}
\newcommand{\pad}{\blue{\bm{\dot{\delta}}}}
\newcommand{\pam}{\blue{\Delta(\pa)}}

\newcommand{\yc}{\blue{y}}
\newcommand{\cge}{\succcurlyeq}
\newcommand{\cle}{\preccurlyeq}
\newcommand{\cl}{\prec}
\newcommand{\cg}{\succ}
\newcommand{\rk}{\text{rank}}
\newcommand{\col}{\text{col}}

%\newcommand{\pa}{\blue{\delta}}
%\newcommand{\setpa}{\blue{\bm{\delta}}}

\renewcommand{\c}[1]{ {\cal #1} }
\newcommand{\bm}[1]{\text{\boldmath$ #1 $} }
\newcommand{\dis}{\displaystyle}
\newcommand{\ab}[1]{\langle #1 \rangle}
\newcommand{\ub}[2]{\underbrace{#1}_{#2}}

\newcommand{\diag}{ \operatornamewithlimits{diag} }
\newcommand{\re}{ \operatorname{Re} }
\newcommand{\he}{ \operatorname{He} }
\renewcommand{\t}[1]{ \tilde{\! #1\,} }
\newcommand{\tc}[1]{ \tilde{{\cal #1}} }
% \renewcommand{\h}[1]{ \hat{#1} } % 2014.08.06: remove this short-hand command, because it cause conflict with Vietnamese language
% If \h is used, some characters with accent like ể will cause error for amsmath: ``\hat allowed only in math mode''
\newcommand{\hc}[1]{ \hat{{\cal #1}} }
\newcommand{\esup}{ \operatornamewithlimits{ess\,sup} }
\newcommand{\tr}{ \operatorname{trace} }
\newcommand{\co}{ \mbox{co} }
\renewcommand{\i}{i}
\newcommand{\io}{i\omega}
\newcommand{\om}{\omega}
\newcommand{\la}{\lambda}
\newcommand{\ga}{\gamma}

\newcommand{\Cl}{ {\mathbb{C}^-} }
\newcommand{\Ci}{ {\mathbb{C}^0} }
\newcommand{\Cr}{ {\mathbb{C}^+} }
\newcommand{\C}{ {\mathbb{C}} }
\newcommand{\Cset}{ {\mathbb{C}} }

\newcommand{\Q}{ {\mathbb{Q}} }
\newcommand{\R}{ {\mathbb{R}} }
\newcommand{\Rset}{ {\mathbb{R}} }
\renewcommand{\H}{ {\mathbb{H}} }
\renewcommand{\S}{ {\mathbb{S}} }
\newcommand{\Nset}{ {\mathbb{N}} }
\newcommand{\Zset}{ {\mathbb{Z}} }
\newcommand{\alert}[1]{{\bf #1}}
%%%%%%%%%%%%%%%%%%%%%%%%%%%%%%%%%%%%%%%%%%%%%%%%%%%%%%%%%%%%%%%%%%%%%%%%%%%%%%%%%%%%%%%%%%%%%%%%%%%%%%%%%%%%%%%%%%%

\newcommand{\Alf}{{D}}
\newcommand{\Blf}{{C}}
\newcommand{\Clf}{{B}}
\newcommand{\Dlf}{{A}}


\newcommand{\Ac}{\green{A_c}}
\newcommand{\Bc}{\green{B_c}}
\newcommand{\Cc}{\green{C_c}}
\newcommand{\Dc}{\green{D_c}}

\newcommand{\Acl}{{\cal A}}
\newcommand{\Bcl}{{\cal B}}
\newcommand{\Ccl}{{\cal C}}
\newcommand{\Dcl}{{\cal D}}
\newcommand{\Xcl}{{ \green{\cal X} }}
\newcommand{\Ycl}{{\cal Y}}
\newcommand{\Zcl}{{\cal Z}}
\newcommand{\Pcl}{{\cal P}}
%closed-loop transformed
\newcommand{\A}{\bm{A}}
\newcommand{\B}{\bm{B}}
%\renewcommand{\C}{\bm{C}}
\newcommand{\D}{\bm{D}}
\renewcommand{\X}{\bm{X}}
\newcommand{\DX}{\bm{Y}}

\renewcommand{\v}{\red{v}}
\newcommand{\K}{\red{K}}
\renewcommand{\L}{\red{L}}
\newcommand{\M}{\red{M}}
\newcommand{\N}{\red{N}}
\renewcommand{\X}{\red{X}}
\newcommand{\Y}{\red{Y}}
\newcommand{\Z}{\red{Z}}
\newcommand{\F}{\blue{F}}

\newcommand{\T}{\c{T}}

\newcommand{\Ri}{\R\cup\{\infty\}}

%%%% Dang added %%%%
\newcommand{\xb}{\mathbf{x}}
\newcommand{\grad}{\nabla}
\newtheorem{Lem}{Bổ đề}
\newtheorem{Rem}{Lưu ý}
\newtheorem{Def}{Định nghĩa}
\newtheorem{Thm}{Định lý}
\newtheorem{Exa}{Ví dụ}



%% end of slides.tex
% % \renewcommand{\pause}{}
% \newcommand{\Hi}{H_\infty}

\begin{document}
% \def\firstofone#1{#1}
% \let\uppercase\firstofone
% \let\MakeUppercase\firstofone

\sdd

% \newpage\titlogo{Tiêu đề có logo lớn}

% \newpage\titf{Tiêu đề có logo khá lớn, dùng làm trang đầu}

% \newpage\titls{Tiêu đề có logo nhỏ}

% spa{1} Tạo một khoảng cách giữa 2 dòng cho dễ nhìn

\newpage\titf{Chương 2: Giải phương trình phi tuyến}
\Large{\textbf{Tính gần đúng nghiệm thực của một phương trình}}

\enu{
\item Đặt bài toán
\item Nghiệm và khoảng cách ly nghiệm
\item Phương pháp chia đôi
\item Phương pháp lặp đơn
\item Phương pháp dây cung
\item Phương pháp Newton
\item Bài tập
% \item Lập trình giải thuật
}
\spa{0}

% \sagetexpause
\newpage\titls{Đặt bài toán}

Tìm nghiệm của phương trình phi tuyến:
\begin{align}\label{eq_baitoan}
 f(x) = 0
\end{align}
với $f$ là hàm liên tục trên một khoảng đóng hoặc mở thuộc $\Rset$.

Vấn đề: khó tính được nghiệm chính xác của \eqref{eq_baitoan}, nên ta cần \textbf{tìm nghiệm gần đúng}.

\begin{tabular}{lr}
Ví dụ:
\begin{sagesilent} 
f(x) = sin(2*x) + x^3 - 2
\end{sagesilent} 
  \begin{tabular}{c}
% \begin{align}\label{eq_ptvidu1}
    $f(x) = \sage{f(x)}$
% \end{align}
  \end{tabular}
&
  \begin{tabular}{c}
    \sageplot[scale=.3]{plot(f, 0, 2)}
  \end{tabular}
\end{tabular}

\spa{0}

\newpage\titls{Đặt bài toán}

Mục tiêu của chương này:

\itemi{
\item Tìm hiểu về cách \textbf{đánh giá sai số} của một nghiệm gần đúng của phương trình phi tuyến một ẩn
\item Học các nguyên lý xác định nghiệm qua \textbf{ý nghĩa hình học} của các giải thuật
\item Biết cách thiết lập các \textbf{công thức tính toán} trong các giải thuật
\item Biết \textbf{áp dụng} các giải thuật (tính toán bằng tay, và lập trình các chương trình chứa giải thuật)
}

\newpage\titls{Nghiệm và khoảng cách ly nghiệm}

Nghiệm chính xác của phương trình $f(x)=\sage{f(x)}=0$ là:
\begin{sagesilent}
s=solve(f,x)
\end{sagesilent}
\[
\sage{s[0]}
\]
\[
\sage{s[1]}
\]
\[
\red{\sage{s[2]}}
\]

\red{(hai giá trị đầu là nghiệm phức, chỉ có nghiệm thứ ba là số thực)}
\spa{1}

Để dễ tìm giá trị số, trước tiên ta tìm \textbf{khoảng cách ly nghiệm}, tức là các khoảng đóng $[a,b]$ (hoặc khoảng mở $(a,b)$) mà trên đó tồn tại duy nhất nghiệm. Sau đó tìm nghiệm gần đúng trong đó.

\textbf{Lưu ý}: Khoảng cách ly nghiệm càng nhỏ thì sai số của nghiệm gần đúng là càng nhỏ, và càng ít tốn công tìm kiếm.

\newpage\titls{Nghiệm và khoảng cách ly nghiệm}

Ví dụ:

\green{Khoảng cách ly nghiệm là $[0.8,1.3]$}

\red{Khoảng cách ly nghiệm là $[1,1.1]$}

\cen{\sageplot[scale=.35]{plot(f, 0, 2)+line([(0.8,-1),(0.8,4)], color='green')+line([(1.3,-1),(1.3,4)], color='green')+line([(1,-1),(1,4)], color='red')+line([(1.1,-1),(1.1,4)], color='red')}}

\newpage\titls{Nghiệm và khoảng cách ly nghiệm}

\begin{Thm}\label{thm_hamtraidau}
 Nếu hàm $f(x)$ liên tục trên đoạn $[a,b]$ và giá trị của hàm trái dấu tại hai đầu mút thì phương trình \eqref{eq_baitoan} có nghiệm trong khoảng $[a,b]$. Thêm nữa, nếu $f(x)$ đơn điệu trên $[a,b]$ thì nghiệm là duy nhất.
\end{Thm}

% \begin{Exa} ×
Ví dụ: Tìm các khoảng cách ly nghiệm của phương trình $f(x)=0$ với
 \begin{sagesilent}
  f(x) = x^3-3*x+1
 \end{sagesilent}

 \begin{tabular}{lc}
    \begin{tabular}{l}
    Hàm số: $f(x)=\sage{f(x)}$\\
    Tính giá trị hàm số ở vài điểm \\
    $\begin{array}{c|rrrrr}
    x & -2 & -1 & 0 & 1 & 2 \\
    \hline
    f(x) & \sage{f(-2)} & \sage{f(-1)} & \sage{f(0)} & \sage{f(1)} & \sage{f(2)}
    \end{array}$\\
    ta tìm được 3 khoảng cách ly nghiệm: \\
    $[-2,-1]$, $[0,1]$, và $[1,2]$
    \end{tabular}
 & 
    \begin{tabular}{c}
    Đồ thị:\\
 \sageplot[scale=.25]{plot(f,-2,2)}
    \end{tabular} 
\end{tabular}
% \end{Exa}

\newpage\titls{Nghiệm và khoảng cách ly nghiệm}

\textbf{Đánh giá sai số:} ta dựa vào \textbf{định lý giá trị trung bình} của Lagrange:

\begin{Thm}\label{thm_gttb}
 Nếu hàm số $f(x)$ liên tục trên $[a,b]$ và có đạo hàm trong $(a,b)$, thì tồn tại một giá trị $c \in (a,b)$ thỏa mãn:
 \begin{equation}
  f'(c) = \frac{f(b)-f(a)}{b-a}
 \end{equation}
\end{Thm}
 \begin{tabular}{lc}
    \begin{tabular}{l}
    \textbf{Ý nghĩa hình học}:\\
    Tiếp tuyến với đường cong $y=f(x)$ \\
    tại điểm $(c,f(c))$ song song với ab
    \end{tabular}
  &
    \begin{tabular}{c}
    \includegraphics[scale=.2]{Lagrange}
    \end{tabular}
 \end{tabular}

\newpage\titls{Nghiệm và khoảng cách ly nghiệm} 

\textbf{Áp dụng Định lý \ref{thm_gttb}}: gọi $p$ là nghiệm chính xác và $x^*$ là nghiệm gần đúng của $f(x)$ trên $[a,b]$. Xét trường hợp $f(x)$ có đạo hàm trên $(p,x^*)$ (nếu $p<x^*$) hoặc trên $(x^*,p)$ (nếu $x^*<p$). Theo định lý Lagrange, tồn tại một điểm $c \in (p,x^*)$ (hoặc $c \in (x^*,p)$), sao cho:
\begin{align*}
 f'(c)&=\frac{f(x^*)-f(p)}{x^*-p} \\
 \Rightarrow |x^*-p|&=\frac{|f(x^*)-f(p)|}{|f'(c)|}=\frac{|f(x^*)|}{|f'(c)|}
\end{align*}

Nếu có số $m$ thỏa: $\mathbf{0<m\leq |f'(x)|, \forall x \in [a,b]}$, thì $|f'(c)| \geq m$, do đó:
$|x^*-p|\leq \frac{|f(x^*)|}{m}$
\bigskip

Vậy ta có thể chọn sai số tuyệt đối là: $\mathbf{\Delta_{x^*}=\frac{|f(x^*)|}{m}}$.

\newpage\titls{Nghiệm và khoảng cách ly nghiệm} 

\textbf{Công thức đánh giá sai số tổng quát:}

\begin{Thm}
 Giả sử hàm $f(x)$ liên tục trên đoạn $[a,b]$, khả vi trong $(a,b)$.  Nếu $x^*$ là nghiệm xấp xỉ của nghiệm chính xác $p \in [a,b]$ và $|f'(x)| \geq m > 0, \forall x \in [a,b]$, thì sai số của $x^*$ thỏa công thức sau:
 \begin{equation}\label{eq_saisotongquat}
  |x^* - p| \leq \frac{|f(x^*)|}{m}
 \end{equation}
\end{Thm}

Ví dụ: xét phương trình $f(x) = \sage{f(x)} = 0$ trong $[1.3,2]$, giả sử tìm được nghiệm gần đúng $x^* = 1.5$. Xét đạo hàm:
\begin{sagesilent}
 f_d1=diff(f, x, 1)(x)
 m = abs(f_d1(1.3))
\end{sagesilent}
\begin{align}
|f'(x)| = |\sage{f_d1}| \geq \sage{round(m,3)} = m, \forall x \in [1.3,2]
\end{align}
Do đó: $|x^* - p| \leq \frac{|f(1.3)|}{|\sage{round(m,3)}|} \approx \sage{round(abs(f(1.3)/m),3)}$  
% RealField(10)(m) to convert the variable to 10-bit binary precision
% m.n(10) also convert the variable to 10-bit binary precision
% round(m,3) to round up to 3 decimal digits

\newpage\titls{Phương pháp chia đôi}

\blue{Giải thuật:
\itemi{
\item Xét đoạn $[a,b]$ là khoảng cách ly nghiệm và $f(a)f(b)<0$. Đặt $a_0=a, b_0=b$ và $x_0=\frac{b+a}{2}$ (điểm giữa của đoạn $[a,b]$).
\item Làm phép lặp (với $k$ là chỉ số vòng lặp, khởi đầu $k=0$): 
\itemi{
  \item Nếu $f(x_k)=0$ thì dừng lại, xác định $x_k$ là nghiệm.
  \item Nếu $\mathbf{f(x_k)f(a_k)}<0$ (tức là $\mathbf{f(a_k)}$ trái dấu với $f(x_k)$) thì đặt $\mathbf{a_{k+1}=a_k}, b_{k+1}=x_k, x_{k+1}=\frac{b_{k+1}+a_{k+1}}{2}$.
  \item Nếu $\mathbf{f(x_k)f(b_k)}<0$ (tức là $\mathbf{f(b_k)}$ trái dấu với $f(x_k)$), khi đó đặt $a_{k+1}=x_k, \mathbf{b_{k+1}=b_k}, x_{k+1}=\frac{b_{k+1}+a_{k+1}}{2}$.
  \item Tăng $k$ thêm 1: $k=k+1$, và lặp lại cách tính toán này.
}
}
}

\newpage\titls{Phương pháp chia đôi}

Kết quả từ giải thuật chia đôi:
\itemi{
\item Ta luôn duy trì được $f(a_k)f(b_k)<0, \forall k$ (trừ khi đạt được nghiệm chính xác). Như vậy $[a_k,b_k]$ luôn là khoảng cách ly nghiệm.
\item Sau mỗi bước, khoảng cách ly nghiệm giảm đi một nửa: $b_{k}-a_{k} = \frac{b_{k-1}-a_{k-1}}{2}$.%dùng k và k-1 để sử dụng ngay vào công thức dưới này, cho SV dễ nhận ra
\item Đặt $d_k=b_k-a_k$, ta có công thức:
\begin{align}\label{eq_khoangcachlyk}
 d_k = \frac{b_{k-1}-a_{k-1}}{2} = \frac{b_{k-2}-a_{k-2}}{2^2} = \cdots = \frac{b_{0}-a_{0}}{2^k} = \frac{b-a}{2^k}
\end{align}
dễ thấy $\lim_{k \to \infty} d_k = 0$, tức là hai dãy số $\{a_k\}, \{b_k\}$ hội tụ. Do $[a_k,b_k]$ là khoảng cách ly nghiệm nên nếu gọi $p$ là nghiệm thì ta luôn có $a_k \leq p \leq b_k$, nên các dãy số $\{a_k\}, \{b_k\}$ hội tụ về $p$.
}

\newpage\titls{Phương pháp chia đôi}

Đánh giá sai số: cho giải thuật chia đôi được lặp cho đến khi đạt được $x_k$ (tức là điểm giữa của đoạn $[a_k,b_k]$). Ta dùng công thức sai số là:
\blue{
\begin{align}
 |x_k - p| \leq \frac{b-a}{2^{k+1}}
\end{align}
}
(Lưu ý: tại sao ở mẫu số có số mũ $k+1$? Vì ta có $d_k= \frac{b-a}{2^{k}}$ theo \eqref{eq_khoangcachlyk}, điểm $x_k$ chia đôi khoảng $[a_k,b_k]$ nên sai số được chia cho 2 nữa)

\newpage\titls{Phương pháp chia đôi}

\begin{sagesilent}
f = sin(2*x) + x^3 - 2
\end{sagesilent}
Ví dụ: giải bài toán $f(x)=\sage{f(x)}=0$ bằng phương pháp chia đôi trên đoạn $[0,1.5]$

\begin{sagesilent}
 xk=range(5)
 ak=range(5)
 bk=range(5)
 fxk=range(5)
 p=find_root(f,0,1.5) # find nearly exact root
 ak[0]=0; bk[0]=1.5; xk[0]=(ak[0]+bk[0])/2; fxk[0]=f(x=xk[0])
 for i in range(4):
    if N(f(x=xk[i])*f(x=ak[i]))<= 0:
        ak[i+1]=ak[i]
        bk[i+1]=xk[i]
    else:    
        ak[i+1]=xk[i]
        bk[i+1]=bk[i]
    xk[i+1]=(bk[i+1]+ak[i+1])/2
    fxk[i+1]=f(x=xk[i+1])
 ak_disp=[round(ak[i],4) for i in range(5)] # to display ak with 4 decimal digits
 bk_disp=[round(bk[i],4) for i in range(5)]
 xk_disp=[round(xk[i],4) for i in range(5)]
 fxk_disp=[round(fxk[i],4) for i in range(5)]
 p_disp=round(p,4)
\end{sagesilent}

Giá trị hàm ở hai đầu: $f(a)=\sage{round(f(x=ak[0]),4)}$, $f(b)=\sage{round(f(x=bk[0]),4)}$

% % hiển thị bảng kết quả bằng cách lập trình Sagetex; cần dùng lệnh latex() của Sage để chuyển giá trị int thành str cho hiển thị latex
\begin{sagesilent}
table=r"\begin{tabular}{ccccc}"
table+=r"$k$ & $a_k$ & $b_k$ & $x_k$ & $f(x_k)$ \\ \hline"
for i in range(5):
  table+=latex(i) + r"&" + latex(ak_disp[i]) + r"&" + latex(bk_disp[i]) + r"&" + latex(xk_disp[i]) + r"&" + latex(fxk_disp[i]) + r"\\"
table+=r"\end{tabular}"
\end{sagesilent}
\cen{\sagestr{table}}

% Để test:  table+=r"0 & 1 & 2 & 3 & 4 \\"


% % Hiển thị các array, kiểu cũ
% \begin{align*}
% \begin{array}{rl}
% a_k: & \sage{ak_disp} \\
% b_k: & \sage{bk_disp} \\
% x_k: & \sage{xk_disp} \\
% f(x_k): & \sage{fxk_disp}
% \end{array}
% \end{align*}


So sánh với nghiệm chính xác: $p = \sage{p_disp}$.


% \begin{sagesilent}
%  a=matrix([[1, 2 , 3],[2, 4, 5],[6,7,8]])
% \end{sagesilent}
% 
% Ma tran:
% $\sage{a}$

\newpage\titls{Phương pháp chia đôi}

Minh họa các bước giải qua đồ thị:

% \cen{\sageplot[scale=.35]{plot(f, -0.5, 2)+line([(ak[0],-1),(ak[0],4)], color='green')+line([(ak[1],-1),(ak[1],4)], color='green')+line([(ak[2],-1),(ak[2],4)], color='green')+line([(ak[3],-1),(ak[3],4)], color='green')+line([(ak[4],-1),(ak[4],4)], color='green')+line([(bk[0],-1),(bk[0],4)], color='blue')+line([(bk[1],-1),(bk[1],4)], color='blue')+line([(bk[2],-1),(bk[2],4)], color='blue')+line([(bk[3],-1),(bk[3],4)], color='blue')+line([(bk[4],-1),(bk[4],4)], color='blue')+line([(xk[0],-1),(xk[0],4)], color='red')+line([(xk[1],-1),(xk[1],4)], color='red')+line([(xk[2],-1),(xk[2],4)], color='red')+line([(xk[3],-1),(xk[3],4)], color='red')+line([(xk[4],-1),(xk[4],4)], color='red')+line([(p,-1.5),(p,4.5)], color='black')}}

\textbf{Bước 1}:

\begin{sagesilent}
 i=0
 plotchiadoi=plot(f, -0.5, 2)+line([(ak[i],-1),(ak[i],4)], color='green')+line([(bk[i],-1),(bk[i],4)], color='blue')+line([(xk[i],-1),(xk[i],4)], color='red')+line([(p,-1.5),(p,4.5)], color='black')
\end{sagesilent}

\cen{\sageplot[scale=.35]{plotchiadoi}}

\newpage\titls{Phương pháp chia đôi}

Minh họa các bước giải qua đồ thị:

\textbf{Bước 2}:

\begin{sagesilent}
 i=1
 plotchiadoi=plot(f, -0.5, 2)+line([(ak[i],-1),(ak[i],4)], color='green')+line([(bk[i],-1),(bk[i],4)], color='blue')+line([(xk[i],-1),(xk[i],4)], color='red')+line([(p,-1.5),(p,4.5)], color='black')
\end{sagesilent}

\cen{\sageplot[scale=.35]{plotchiadoi}}

\newpage\titls{Phương pháp chia đôi}

Minh họa các bước giải qua đồ thị:

\textbf{Bước 3}:

\begin{sagesilent}
 i=2
 plotchiadoi=plot(f, -0.5, 2)+line([(ak[i],-1),(ak[i],4)], color='green')+line([(bk[i],-1),(bk[i],4)], color='blue')+line([(xk[i],-1),(xk[i],4)], color='red')+line([(p,-1.5),(p,4.5)], color='black')
\end{sagesilent}

\cen{\sageplot[scale=.35]{plotchiadoi}}

\newpage\titls{Phương pháp chia đôi}

Minh họa các bước giải qua đồ thị:

\textbf{Bước 4}:

\begin{sagesilent}
 i=3
 plotchiadoi=plot(f, -0.5, 2)+line([(ak[i],-1),(ak[i],4)], color='green')+line([(bk[i],-1),(bk[i],4)], color='blue')+line([(xk[i],-1),(xk[i],4)], color='red')+line([(p,-1.5),(p,4.5)], color='black')
\end{sagesilent}

\cen{\sageplot[scale=.35]{plotchiadoi}}

\newpage\titls{Phương pháp chia đôi}

Minh họa các bước giải qua đồ thị:

\textbf{Bước 5}:

\begin{sagesilent}
 i=4
 plotchiadoi=plot(f, -0.5, 2)+line([(ak[i],-1),(ak[i],4)], color='green')+line([(bk[i],-1),(bk[i],4)], color='blue')+line([(xk[i],-1),(xk[i],4)], color='red')+line([(p,-1.5),(p,4.5)], color='black')
\end{sagesilent}

\cen{\sageplot[scale=.35]{plotchiadoi}}

\newpage\titls{Phương pháp chia đôi}

\begin{sagesilent}
f = x^2 - 2
\end{sagesilent}
Ví dụ (\red{tính căn bậc hai của 2}): giải bài toán $f(x)=\sage{f(x)}=0$ bằng phương pháp chia đôi trên đoạn $[0,2]$

\begin{sagesilent}
 n=7
 xk=range(n)
 ak=range(n)
 bk=range(n)
 fxk=range(n)
 p=find_root(f,0,2) # find nearly exact root
 ak[0]=0; bk[0]=2; xk[0]=(ak[0]+bk[0])/2; fxk[0]=f(x=xk[0])
 for i in range(n-1):
    if N(f(x=xk[i])*f(x=ak[i]))<= 0:
        ak[i+1]=ak[i]
        bk[i+1]=xk[i]
    else:    
        ak[i+1]=xk[i]
        bk[i+1]=bk[i]
    xk[i+1]=(bk[i+1]+ak[i+1])/2
    fxk[i+1]=f(x=xk[i+1])
 ak_disp=[round(ak[i],4) for i in range(n)] # to display ak with 4 decimal digits
 bk_disp=[round(bk[i],4) for i in range(n)]
 xk_disp=[round(xk[i],4) for i in range(n)]
 fxk_disp=[round(fxk[i],4) for i in range(n)]
 p_disp=round(p,4)
\end{sagesilent}

Giá trị hàm ở hai đầu: $f(a)=\sage{round(f(x=ak[0]),4)}$, $f(b)=\sage{round(f(x=bk[0]),4)}$

% % hiển thị bảng kết quả bằng cách lập trình Sagetex; cần dùng lệnh latex() của Sage để chuyển giá trị int thành str cho hiển thị latex
\begin{sagesilent}
table=r"\begin{tabular}{ccccc}"
table+=r"$k$ & $a_k$ & $b_k$ & $x_k$ & $f(x_k)$ \\ \hline"
for i in range(n):
  table+=latex(i) + r"&" + latex(ak_disp[i]) + r"&" + latex(bk_disp[i]) + r"&" + latex(xk_disp[i]) + r"&" + latex(fxk_disp[i]) + r"\\"
table+=r"\end{tabular}"
\end{sagesilent}
\cen{\sagestr{table}}

Nghiệm chính xác: $p=\sage{p_disp}$
% \sagetexunpause

\newpage\titls{Phương pháp lặp đơn}

Đầu tiên, chuyển từ phương trình \eqref{eq_baitoan} sang dạng tương đương:
\begin{align}\label{eq_baitoanlap}
 x = g(x)
\end{align}

Nghiệm của phương trình \eqref{eq_baitoanlap} được gọi là \textbf{điểm bất động} của hàm số $g(x)$.

\blue{Giải thuật:
\itemi{
\item Xét đoạn $[a,b]$ là khoảng cách ly nghiệm. Chọn giá trị đầu $x_0 \in [a,b]$.
\item Lặp: $x_k=g(x_{k-1})$, với $k=1, 2, ...$
}
}

Ta sẽ khảo sát phương pháp này có \textbf{hội tụ} về nghiệm của bài toán hay không, tùy vào đặc điểm của hàm $g(x)$.

\newpage\titls{Phương pháp lặp đơn}

Ý nghĩa hình học:
\itemi{
\item Tìm điểm giao của đường thẳng $y=x$ và đồ thị $y=g(x)$.
\item Giải thuật lặp: khi có giá trị $x_{k-1}$, ta lấy điểm đối xứng của điểm $g(x_{k-1})$ qua đường $y=x$, chiếu xuống trục hoành được $x_k$.
}

%  Đây là hàm chuyển đổi cho f(x)=x^3-3*x+1 đã có ở phía trên
\begin{sagesilent}
 g1 = x^2/2 - sin(x)
 g1_diff=diff(g1, x, 1)
 x_k=1
 x_kplus1=g1(x=x_k)
 plotg1=plot(g1, -1, 2, color='red')
 plotyeqx=plot(x, -1, 2, color='blue')
 plotxkchieu=line([(x_k,x_kplus1),(x_kplus1,x_k)], color='green', linestyle='-')
 plotxkline_ngang=line([(0,x_kplus1),(x_k,x_kplus1)], color='green', linestyle='--')
 plotxkline_doc=line([(x_kplus1,x_k),(x_kplus1,0)], color='green', linestyle='--') 
 cxk=circle((x_k,g1(x=x_k)),0.04,color='green')
 cxkplus1=circle((x_kplus1,g1(x=x_kplus1)),0.04,color='red')
\end{sagesilent}

 \begin{tabular}{lc}
    \begin{tabular}{l}
    \red{Hàm số: $g(x)=\sage{g1}$}\\
    \blue{Đường $y=x$}\\
    \green{Điểm bắt đầu: $x_{k-1}$}\\
    \red{Điểm tiếp theo: $x_k$}
    \end{tabular}
  &
    \begin{tabular}{c}
    \sageplot[scale=.3]{plotg1+plotyeqx+plotxkchieu+plotxkline_ngang+plotxkline_doc+cxk+cxkplus1}
    \end{tabular}
 \end{tabular}


\newpage\titls{Phương pháp lặp đơn}

Định nghĩa: Hàm $g(x)$ được gọi là \textbf{hàm co} trong đoạn $[a,b]$ nếu tồn tại một số $q: 0 \leq q < 1$ sao cho:
\begin{equation}
 |g(x_1) - g(x_2)| \leq q |x_1 - x_2|, \forall x_1, x_2 \in [a,b]
\end{equation}
Giá trị $q$ được gọi là \textbf{hệ số co}.

\begin{Thm}
 Nếu $g(x)$ là hàm co trên đoạn $[a,b]$, thì nó liên tục trên đó.
\end{Thm}

\begin{Thm}
 Nếu $g(x)$ liên tục trên đoạn $[a,b]$, khả vi trong $(a,b)$ và $\exists q: 0 \leq q < 1$ sao cho $|g'(x)| \leq q, \forall x \in (a,b)$, thì $g(x)$ là hàm co trên $[a,b]$ với hệ số co là $q$.
\end{Thm}

\newpage\titls{Phương pháp lặp đơn}

\begin{sagesilent}
 f2(x) = x^3 + x - 10
 g2 = (10 - x)^(1/3)
 g2_diff=diff(g2, x, 1)
\end{sagesilent}

Ví dụ: Xét phương trình $f(x)=0$ với $f(x)=\sage{f2(x)}$ trên đoạn $[0,3]$, chuyển sang dạng \eqref{eq_baitoanlap}, ta được: $g(x) = \sage{g2}$.

Ta có:
\begin{align*}
 |g'(x)| = \Bigg| \sage{g2_diff} \Bigg| \leq \sage{abs(g2_diff(x=3))} \approx \sage{round(abs(g2_diff(x=3)),4)} \triangleq q < 1
\end{align*}
Do đó $g(x)$ là hàm co trên đoạn $[0,3]$.

% Ví dụ khác: Xét $f(x)=\sage{f}$ trên đoạn $[0,1.5]$:
% \begin{sagesilent}
%  g2 = (10 - x)^(1/3)
%  g2_diff1=diff(g2, x, 1)
% \end{sagesilent}

\newpage\titls{Phương pháp lặp đơn}

Định lý sau đây tạo cơ sở cho phương pháp lặp đơn, còn gọi là \textbf{nguyên lý ánh xạ co}.
\begin{Thm}
 Giả sử $g(x)$ là hàm co trên đoạn $[a,b]$ với hệ số co là $q$; đồng thời $\forall x \in [a,b], g(x) \in [a,b]$. Khi đó, với mọi giá trị ban đầu $x_0 \in [a,b]$, dãy lặp $\{x_n\}_{n=1}^\infty$ của giải thuật lặp sẽ hội tụ về nghiệm $p$ của phương trình \eqref{eq_baitoanlap}.
\end{Thm}

Sai số thu được:
\begin{align}
 |x_n - p| \leq \frac{q^n}{1-q} |x_1 - x_0|
\end{align}
hoặc:
\begin{align}
 |x_n - p| \leq \frac{q}{1-q} |x_n - x_{n-1}|
\end{align}

\newpage\titls{Phương pháp lặp đơn}

Tóm tắt tính chất hội tụ:
\itemi{
\item Hàm số là \emph{co} (hệ số co thỏa mãn $q<1$)
\item Tập hợp $[a,b]$ là \emph{tập giới nội} của hàm số $g(x)$
\item Tốc độ hội tụ phụ thuộc vào giá trị của hệ số co $q$: $q$ càng nhỏ thì hội tụ càng nhanh (ít bước).
}

\newpage\titls{Phương pháp lặp đơn}

Ví dụ: áp dụng phương pháp lặp để giải $x=g(x)=\sage{g2}$ với $x_0=0$

\cen{\sageplot[scale=.4]{plot(g2,-1,10)+line([(-1,-1),(6,6)],color='green')}}

\newpage\titls{Phương pháp lặp đơn}

Ví dụ: áp dụng phương pháp lặp để giải $x=g(x)=\sage{g2}$ với $x_0=0$

\begin{sagesilent}
 n=6
 decipoint=4
 xk=range(n+1)
 g2xk=range(n)
 deltaxk=range(n)
 xk[0]=0
 p_g2=find_root(f2,0,5) # find nearly exact root
 for i in range(n):
    g2xk[i]=g2(x=xk[i])
    xk[i+1]=g2xk[i]
    deltaxk[i]=abs(xk[i]-p_g2)
 xk_disp=[round(xk[i],decipoint) for i in range(n)]
 g2xk_disp=[round(g2xk[i],decipoint) for i in range(n)]
 deltaxk_disp=[round(deltaxk[i],decipoint) for i in range(n)]
 p_g2_disp=round(p_g2,decipoint)
 table2=r"\begin{tabular}{c|ccc}"
 table2+=r"$k$ & $x_k$ & $g(x_k)$ & $|x_k-p|$ \\ \hline"
 for i in range(n):
   table2+=latex(i) + r"&" + latex(xk_disp[i]) + r"&" + latex(g2xk_disp[i]) + r"&" + latex(deltaxk_disp[i]) + r"\\"
 table2+=r"\end{tabular}"
\end{sagesilent}
% Nếu decipoint=5, khi tính toán sẽ cho ra sai số kiểu 9x10^(-5), khi compile Latex sẽ có báo lỗi kiểu "Missing $ inserted \cen{\sagestr{table2}}", "Extra }, or forgotten $ \cen{\sagestr{table2}}", chắc là do Sagetex có đưa ra ký tự x là đặc biệt. Dù sao khi đó file PDF vẫn compile được, có thể bỏ qua lỗi đó. Nếu decipoint=4 (làm tròn thêm nữa) thì không có báo lỗi.
\cen{\sagestr{table2}}

So sánh với nghiệm chính xác: $p=\sage{p_g2_disp}$


\newpage\titls{Phương pháp lặp đơn}

\begin{sagesilent}
 g3=cos(x)
\end{sagesilent}

Ví dụ: áp dụng phương pháp lặp để giải $x=g(x)=\sage{g3}$ với $x_0=0$

\cen{\sageplot[scale=.4]{plot(g3,-pi/2,pi/2)+line([(-1,-1),(1,1)],color='green')}}

\newpage\titls{Phương pháp lặp đơn}

Ví dụ: áp dụng phương pháp lặp để giải $x=g(x)=\sage{g3}$ với $x_0=0$

\begin{sagesilent}
 n=9
 decipoint=5
 xk=range(n+1)
 g3xk=range(n)
 deltaxk=range(n)
 xk[0]=0
 p_g3=find_root(g3-x,0,1) # find nearly exact root
 for i in range(n):
    g3xk[i]=g3(x=xk[i])
    xk[i+1]=g3xk[i]
    deltaxk[i]=abs(xk[i]-p_g3)
 xk_disp=[round(xk[i],decipoint) for i in range(n)]
 g3xk_disp=[round(g3xk[i],decipoint) for i in range(n)]
 deltaxk_disp=[round(deltaxk[i],decipoint) for i in range(n)]
 p_g3_disp=round(p_g3,decipoint)
 table3=r"\begin{tabular}{c|ccc}"
 table3+=r"$k$ & $x_k$ & $g(x_k)$ & $|x_k-p|$ \\ \hline"
 for i in range(n):
   table3+=latex(i) + r"&" + latex(xk_disp[i]) + r"&" + latex(g3xk_disp[i]) + r"&" + latex(deltaxk_disp[i]) + r"\\"
 table3+=r"\end{tabular}"
\end{sagesilent}

\cen{\sagestr{table3}}

So sánh với nghiệm chính xác: $p=\sage{p_g3_disp}$


% \sagetexpause

\newpage\titls{Phương pháp dây cung}

Xét bài toán \eqref{eq_baitoan} trong khoảng cách ly nghiệm $[a,b]$, hàm $f(x)$ là liên tục và $f(a).f(b)<0$. Ý nghĩa: hai đầu hàm số ở hai bên trục hoành.

\blue{Giải thuật \textbf{lặp} (ở bước thứ $k$, với $a_0=a, b_0=b$):
\itemi{
\item Kẻ đoạn thẳng (\textit{dây cung}) nối hai điểm $(a_k, f(a_k))$ và $(b_k, f(b_k))$ (hai đầu của đồ thị hàm số trong đoạn $[a_k,b_k]$).
\item Dây cung giao với trục hoành tại một điểm, thì lấy hoành độ của điểm đó làm $x_k$.
\item Lặp lại việc dựng dây cung với đoạn nhỏ hơn: Nếu $f(a_k).f(x_k)<0$, thì đặt $a_{k+1}=a_k, \mathbf{b_{k+1}=x_k}$; nếu $f(x_k).f(b)<0$, thì đặt $\mathbf{a_{k+1}=x_k}, b_{k+1}=b_k$.
}
}

\newpage\titls{Phương pháp dây cung}

\begin{sagesilent}
 f = (x-1)^3*e^(x+1)/10-1
\end{sagesilent}
% Sageplot xuat ra hinh khong tot neu nhan 10 ham f len: ti le o truc tung va truc hoanh khong tu dong scale cho phu hop voi khung hinh xuat ra, lam cho hinh rat la dai
% Để sửa tỉ lệ, có thể dùng thuộc tính: p.set_aspect_ratio(r), với p là một đối tượng plot, tỉ lệ r=độ dài trên trục tung/độ dài trên trục hoành

Minh họa: \green{dây cung} của \blue{$f(x)=\sage{f}$} trên đoạn $[1,2]$
% Tọa độ của điểm giao:
% Dây cung: y=f(a)+[(f(b)-f(a))/(b-a)]*(x-a)
% Giao với trục hoành: y=0 <=> x=a-f(a)*(b-a)/(f(b)-f(a))
\begin{sagesilent}
 a=1
 b=2
 plot_hamf=plot(f,a,b, color='blue')
 plot_daycung=line([(a,f(x=a)),(b,f(x=b))], color='green')
 x_k_daycung=a-f(x=a)*(b-a)/(f(x=b)-f(x=a))
 c_xk_daycung=circle((x_k_daycung,0),0.04, color='red')
 c_fxk_daycung=circle((x_k_daycung,f(x=x_k_daycung)),0.04, color='blue')
 plot_daycung2=line([(x_k_daycung,f(x=x_k_daycung)),(b,f(x=b))], color='purple')
\end{sagesilent}

\red{$x=\sage{round(x_k_daycung,4)}$}, \blue{$f(x)=\sage{round(f(x=x_k_daycung),4)}$}, ta có $f(x).f(b)<0$, vậy tiếp theo xét trên đoạn $[\sage{round(x_k_daycung,4)},\sage{b}]$.

\cen{\sageplot[scale=.35]{plot_hamf+plot_daycung+c_xk_daycung+c_fxk_daycung+plot_daycung2}}

% \newpage\titls{Phương pháp dây cung}
% 
% Có 4 trường hợp xảy ra của thuật toán:
% \enu{
% \item $f(a)<0, f(b)>0$, và $f(x_k)<0$: khoảng cách ly nghiệm tiếp theo sẽ là $[x_k, b]$
% \item $f(a)<0, f(b)>0$, và $f(x_k)>0$
% \item $f(a)>0, f(b)<0$, và $f(x_k)<0$
% \item $f(a)>0, f(b)<0$, và $f(x_k)>0$
% }

\newpage\titls{Phương pháp dây cung}

\textbf{Tìm công thức tính giá trị của $x_k$:}

Phương trình của dây cung: 
\begin{equation}\label{eq_daycung}
y=f(a)+\frac{f(b)-f(a)}{b-a}(x-a) 
\end{equation}
Phương trình của trục hoành: $y=0$

Tọa độ của điểm giao giữa dây cung với trục hoành là nghiệm chung của hai phương trình trên. Ta thay $y=0$ vào phương trình \eqref{eq_daycung}:
\begin{align}
y=0 &\Leftrightarrow f(a)+\frac{f(b)-f(a)}{b-a}(x-a)=0 \nonumber\\
&\Leftrightarrow x=a-f(a)\frac{b-a}{f(b)-f(a)} \nonumber\\
&\Leftrightarrow x=a-f(a)\frac{a-b}{f(a)-f(b)} \label{eq_ppdaycung_xtheoa}
\end{align}

\newpage\titls{Phương pháp dây cung}

\textbf{Tìm công thức tính giá trị của $x_k$, cách khác:}

Phương trình của dây cung, \textbf{thiết lập theo $f(b)$}:
\begin{equation}\label{eq_daycung2}
y=f(b)-\frac{f(b)-f(a)}{b-a}(b-x) 
\end{equation}
Phương trình của trục hoành: $y=0$

Tọa độ của điểm giao giữa dây cung với trục hoành là nghiệm chung của hai phương trình trên. Ta thay $y=0$ vào phương trình \eqref{eq_daycung2}:
\begin{align}
y=0 &\Leftrightarrow f(b)-\frac{f(b)-f(a)}{b-a}(b-x)=0 \nonumber\\
&\Leftrightarrow x=b-f(b)\frac{b-a}{f(b)-f(a)} \label{eq_ppdaycung_xtheob}
\end{align}

\textbf{Lưu ý}: cả hai công thức \eqref{eq_ppdaycung_xtheoa} và \eqref{eq_ppdaycung_xtheob} đều đúng.

\newpage\titls{Phương pháp dây cung}

\textbf{Xét mối quan hệ giữa $x_{k+1}$ và $x_k$:}

Theo giải thuật, ta biết rằng hoặc $a_{k+1}=x_k$, hoặc $b_{k+1}=x_k$.% Có hai trường hợp:

\blue{
\itemi{
\item Nếu $a_{k+1}=x_k$, thì dùng công thức \eqref{eq_ppdaycung_xtheoa}, ta được:
\begin{align*}
x_{k+1}&=a_{k+1}-f(a_{k+1})\frac{a_{k+1}-b_{k+1}}{f(a_{k+1})-f(b_{k+1})}\\
\Leftrightarrow x_{k+1}&=x_k-\frac{f(x_k)(x_k-b_{k+1})}{f(x_k)-f(b_{k+1})}
\end{align*}
\item Nếu $b_{k+1}=x_k$, thì dùng công thức \eqref{eq_ppdaycung_xtheob}, ta được:
\begin{align*}
x_{k+1}&=b_{k+1}-f(b_{k+1})\frac{b_{k+1}-a_{k+1}}{f(b_{k+1})-f(a_{k+1})}\\
\Leftrightarrow x_{k+1}&=x_k-\frac{f(x_k)(x_k-a_{k+1})}{f(x_k)-f(a_{k+1})}
\end{align*}
}
}

% \sagetexunpause

\newpage\titls{Phương pháp dây cung}

\begin{sagesilent}
f = sin(2*x) + x^3 - 2
\end{sagesilent}

\begin{sagesilent}
 n=10
 decipoint=4
 xk=range(n)
 ak=range(n)
 bk=range(n)
 deltaxk=range(n)
 fxk=range(n)
 fak=range(n)
 fbk=range(n)
 p_f=find_root(f,0,1.5) # find nearly exact root
 ak[0]=0; bk[0]=1.5;
 for i in range(n-1):
    fak[i]=f(x=ak[i])
    fbk[i]=f(x=bk[i])
    xk[i]=bk[i]-(fbk[i]*(bk[i]-ak[i]))/(fbk[i]-fak[i])
    fxk[i]=f(x=xk[i])
    deltaxk[i]=abs(xk[i]-p_f)
    if N(fxk[i]*fak[i])<= 0:
        ak[i+1]=ak[i]
        bk[i+1]=xk[i]
    else:    
        ak[i+1]=xk[i]
        bk[i+1]=bk[i]
 ak_disp=[round(ak[i],decipoint) for i in range(n)]
 bk_disp=[round(bk[i],decipoint) for i in range(n)]
 xk_disp=[round(xk[i],decipoint) for i in range(n)]
 fxk_disp=[round(fxk[i],decipoint) for i in range(n)]
 deltaxk_disp=[round(deltaxk[i],decipoint) for i in range(n)]
 p_f_disp=round(p_f,decipoint)
 tableDC1=r"\begin{tabular}{c|ccc}"
 tableDC1+=r"$k$ & $x_k$ & $f(x_k)$ & $|x_k-p|$ \\ \hline"
 for i in range(n-1):
   tableDC1+=latex(i) + r"&" + latex(xk_disp[i]) + r"&" + latex(fxk_disp[i]) + r"&" + latex(deltaxk_disp[i]) + r"\\"
 tableDC1+=r"\end{tabular}"
\end{sagesilent}

Ví dụ: giải bài toán $f(x)=\sage{f(x)}=0$ bằng phương pháp dây cung trên đoạn $[0,1.5]$

\cen{\sageplot[scale=.4]{plot(f,0,1.5)}}

\newpage\titls{Phương pháp dây cung}

Ví dụ: giải bài toán $f(x)=\sage{f(x)}=0$ bằng phương pháp dây cung trên đoạn $[0,1.5]$
 (nghiệm chính xác: $p=\sage{p_f_disp}$)

\cen{\sagestr{tableDC1}}

\newpage\titls{Phương pháp dây cung}

\begin{sagesilent}
f = x^3 + x - 10
\end{sagesilent}

\begin{sagesilent}
 n=10
 decipoint=4
 xk=range(n)
 ak=range(n)
 bk=range(n)
 deltaxk=range(n)
 fxk=range(n)
 fak=range(n)
 fbk=range(n)
 p_f=find_root(f,0,3) # find nearly exact root
 ak[0]=0; bk[0]=3;
 for i in range(n-1):
    fak[i]=f(x=ak[i])
    fbk[i]=f(x=bk[i])
    xk[i]=bk[i]-(fbk[i]*(bk[i]-ak[i]))/(fbk[i]-fak[i])
    fxk[i]=f(x=xk[i])
    deltaxk[i]=abs(xk[i]-p_f)
    if N(fxk[i]*fak[i])<= 0:
        ak[i+1]=ak[i]
        bk[i+1]=xk[i]
    else:    
        ak[i+1]=xk[i]
        bk[i+1]=bk[i]
 ak_disp=[round(ak[i],decipoint) for i in range(n)]
 bk_disp=[round(bk[i],decipoint) for i in range(n)]
 xk_disp=[round(xk[i],decipoint) for i in range(n)]
 fxk_disp=[round(fxk[i],decipoint) for i in range(n)]
 deltaxk_disp=[round(deltaxk[i],decipoint) for i in range(n)]
 p_f_disp=round(p_f,decipoint)
 tableDC2=r"\begin{tabular}{c|ccc}"
 tableDC2+=r"$k$ & $x_k$ & $f(x_k)$ & $|x_k-p|$ \\ \hline"
 for i in range(n-1):
   tableDC2+=latex(i) + r"&" + latex(xk_disp[i]) + r"&" + latex(fxk_disp[i]) + r"&" + latex(deltaxk_disp[i]) + r"\\"
 tableDC2+=r"\end{tabular}"
\end{sagesilent}

Ví dụ: giải bài toán $f(x)=\sage{f(x)}=0$ bằng phương pháp dây cung trên đoạn $[0,3]$
 (nghiệm chính xác: $p=\sage{p_f_disp}$)

\cen{\sageplot[scale=.4]{plot(f,-1,4)}}

\newpage\titls{Phương pháp dây cung}

Ví dụ: giải bài toán $f(x)=\sage{f(x)}=0$ bằng phương pháp dây cung trên đoạn $[0,3]$
 (nghiệm chính xác: $p=\sage{p_f_disp}$)

\cen{\sagestr{tableDC2}}


% \sagetexpause

% % Mang về đằng trước, chứng minh công thức sai số tổng quát
% \newpage\titls{Phương pháp dây cung}
% 
% \textbf{Đánh giá sai số của phương pháp dây cung:} ta dựa vào \textbf{định lý giá trị trung bình} của Lagrange:
% 
% \begin{Thm}
%  Nếu hàm số $f(x)$ liên tục trên $[a,b]$ và có đạo hàm trong $(a,b)$, thì tồn tại một giá trị $c \in (a,b)$ thỏa mãn:
%  \begin{equation}
%   f'(c) = \frac{f(b-f(a)}{b-a}
%  \end{equation}
% \end{Thm}
%  \begin{tabular}{lc}
%     \begin{tabular}{l}
%     Ý nghĩa hình học:\\
%     Tiếp tuyến với đường cong $y=f(x)$ \\
%     tại điểm $(c,f(c))$ song song với AB
%     \end{tabular}
%   &
%     \begin{tabular}{c}
%     \includegraphics[scale=.2]{Lagrange}
%     \end{tabular}
%  \end{tabular}
% 
% \newpage\titls{Phương pháp dây cung} 
% 
% \textbf{Đánh giá sai số của phương pháp dây cung:}
% 
% Áp dụng: gọi $p$ là nghiệm của $f(x)$ trên $[a,b]$ (tức là $f(p)=0$). Xét trường hợp $f(x)$ có đạo hàm trên $(p,x_k)$ (nếu $p<x_k$) hoặc trên $(x_k,p)$ (nếu $x_k<p$). Theo định lý Lagrange, tồn tại một điểm $c \in (p,x_k)$ (hoặc $c \in (x_k,p)$), sao cho:
% \begin{align*}
%  f'(c)&=\frac{f(x_k)-f(p)}{x_k-p} \\
%  \Rightarrow |x_k-p|&=\frac{|f(x_k)-f(p)|}{|f'(c)|}=\frac{|f(x_k)|}{|f'(c)|}
% \end{align*}
% 
% Nếu có số $m$ thỏa: $\mathbf{0<m\leq f'(x), \forall x \in [a,b]}$, thì $|x_k-p|\leq \frac{|f(x_k)|}{m}$
% \bigskip
% 
% Vậy ta có thể chọn sai số tuyệt đối là: $\mathbf{\Delta_{x_k}=\frac{|f(x_k)|}{m}}$.

\newpage\titls{Phương pháp dây cung}

Tên giải thuật bằng tiếng Anh:

\itemi{
\item \emph{Secant method}: là giải thuật lặp 
\begin{equation*}
x_{k+1}=x_k - \frac{f(x_k)(x_k-x_{k-1})}{f(x_k)-f(x_{k-1})} 
\end{equation*}
giống như phương pháp dây cung nhưng không so sánh dấu của $f(x)$ với hai đầu $f(a), f(b)$ mà tuần tự bỏ đầu $a_k$ rồi $b_{k+1}$. Giải thuật này có thể hội tụ (tìm được nghiệm) hoặc phân kỳ (không tiến đến nghiệm).
\item \emph{False position method}: là phương pháp dây cung ta học ở đây, luôn đảm bảo có nghiệm, nhờ có kiểm tra dấu của $f(x)$ ở mỗi bước để đảm bảo khoảng tiếp theo vẫn là khoảng cách ly nghiệm.
}

\newpage\titls{Phương pháp Newton}

Còn gọi là \emph{Phương pháp tiếp tuyến} hoặc \emph{Phương pháp Newton-Raphson}.
\bigskip

Xét bài toán \eqref{eq_baitoan} trong khoảng cách ly nghiệm $[a,b]$, với hàm số \textbf{$f(x)$ có đạo hàm} trong $[a,b]$. Ý nghĩa: đồ thị của $f(x)$ là một đường thẳng hoặc đường cong trơn trong $[a,b]$.

\blue{Giải thuật \textbf{lặp} (với điểm ban đầu $x_0 \in [a,b]$):
\itemi{
\item Từ điểm có hoành độ $x_{k-1}$ trên đồ thị của đường cong $f(x)$, ta kẻ tiếp tuyến với đường cong.
\item Đường tiếp tuyến giao với trục hoành tại một điểm, thì lấy hoành độ của điểm đó làm $x_k$.
}
}

\newpage\titls{Phương pháp Newton}

Phương trình của tiếp tuyến là:
\begin{align*}
 y - f(x_{k-1}) = f'(x_{k-1})(x - x_{k-1})
\end{align*}
Điểm giao của đường này với trục hoành là khi $y=0$, tức là:
\blue{
\begin{align}\label{eq_lapNewton}
 \mathbf{x_k = x_{k-1} - \frac{f(x_{k-1})}{f'(x_{k-1})}}, ~ k=1,2,3,...
\end{align}
}

Định lý về sự hội tụ của phép lặp này:
\begin{Thm}
 Giả sử hàm $f(x)$ có đạo hàm đến cấp hai liên tục và các đạo hàm $f'(x)$ và $f''(x)$ không đổi dấu trên đoạn $[a,b]$. Khi đó nếu chọn $x_0$ thỏa điều kiện Fourier: $f(x_0) f''(x_0)>0$, thì dãy lặp $\{x_n\}_{n=1}^\infty$ xác định theo công thức \eqref{eq_lapNewton} sẽ hội tụ về nghiệm của bài toán \eqref{eq_baitoan}.
\end{Thm}

\newpage\titls{Phương pháp Newton}

Lưu ý: 
\itemi{
\item Hệ số góc của tiếp tuyến là giá trị của đạo hàm tại điểm ta vẽ ($f'(x_{k-1})$). Khi hệ số góc bằng 0 tức là đường tiếp tuyến song song với trục hoành (sẽ không tồn tại điểm giao), giải thuật thường được điều chỉnh bằng cách bỏ đi điểm $c$ có $f'(c)=0$ và dùng lại một trong hai nửa đoạn $[a,c]$ hoặc $[c,b]$.
\item Điều kiện Fourier chỉ là điều kiện đủ, không phải là điều kiện cần. Nghĩa là giải thuật cũng có thể hội tụ mà không thỏa điều kiện Fourier.
\item Đánh giá sai số của phương pháp Newton theo công thức sai số tổng quát \eqref{eq_saisotongquat}.
}

% \sagetexunpause
\newpage\titls{Phương pháp Newton}

% \begin{sagesilent}
%  f = e^(x/5) - 5
% \end{sagesilent}
% 
% Minh họa: \green{tiếp tuyến} của \blue{$f(x)=\sage{f}$} với \green{$x_0=10$}
% % Tọa độ của điểm giao:
% % Tiếp tuyến: y - f(x_{k-1}) = f'(x_{k-1})(x - x_{k-1})
% % Giao với trục hoành: y=0 <=> x=[f'(x_{k-1})*x_{k-1}-f(x_{k-1})]/f'(x_{k-1})
% % (vẽ tiếp tuyến trong đoạn [c,b] để khỏi dài quá làm mất cân đối đồ thị
% \begin{sagesilent}
%  a=2.5
%  b=11
%  c=7
%  x0=10
%  fx0=f(x=x0)
%  f_diff=diff(f,x,1)
%  f_diff_x0=f_diff(x=x0)
%  tieptuyen_x0=f_diff_x0*(x - x0)+fx0
%  plot_hamf=plot(f,a,b, color='blue')
%  plot_tieptuyen=plot(tieptuyen_x0,c,b, color='green')
%  x_k_tieptuyen=(f_diff_x0*x0-fx0)/f_diff_x0
%  c_xk_tieptuyen=circle((x_k_tieptuyen,0),0.04, color='red')
%  c_fxk_tieptuyen=circle((x_k_tieptuyen,f(x=x_k_tieptuyen)),0.04, color='blue')
%  c_fx0=circle((x0,fx0),0.04, color='green')
%  plotall_tieptuyen=plot_hamf+plot_tieptuyen+c_xk_tieptuyen+c_fxk_tieptuyen+c_fx0
%  plotall_tieptuyen.set_aspect_ratio(1)
% \end{sagesilent}

\begin{sagesilent}
 f = e^(x) - 5
\end{sagesilent}

Minh họa (\red{tính $ln(5)$}): \green{tiếp tuyến} của \blue{$f(x)=\sage{f}$} với \green{$x_0=2$}
% Tọa độ của điểm giao:
% Tiếp tuyến: y - f(x_{k-1}) = f'(x_{k-1})(x - x_{k-1})
% Giao với trục hoành: y=0 <=> x=[f'(x_{k-1})*x_{k-1}-f(x_{k-1})]/f'(x_{k-1})
% (vẽ tiếp tuyến trong đoạn [c,b] để khỏi dài quá làm mất cân đối đồ thị)
\begin{sagesilent}
 a=1.5
 b=2.2
 c=1.6
 x0=2
 fx0=f(x=x0)
 f_diff=diff(f,x,1)
 f_diff_x0=f_diff(x=x0)
 tieptuyen_x0=f_diff_x0*(x - x0)+fx0
 plot_hamf=plot(f,a,b, color='blue')
 plot_tieptuyen=plot(tieptuyen_x0,c,b, color='green')
 x_k_tieptuyen=(f_diff_x0*x0-fx0)/f_diff_x0
 c_xk_tieptuyen=circle((x_k_tieptuyen,0),0.04, color='red')
 c_fxk_tieptuyen=circle((x_k_tieptuyen,f(x=x_k_tieptuyen)),0.04, color='blue')
 c_fx0=circle((x0,fx0),0.04, color='green')
 plotall_tieptuyen=plot_hamf+plot_tieptuyen+c_xk_tieptuyen+c_fxk_tieptuyen+c_fx0
 plotall_tieptuyen.set_aspect_ratio(0.2)
\end{sagesilent}
% Thuộc tính set_aspect_ratio(r) giúp thay đổi tỉ lệ đồ thị manually: r=tỉ lệ trục tung/trục hoành

\red{$x_1=\sage{round(x_k_tieptuyen,4)}$}, \blue{$f(x_1)=\sage{round(f(x=x_k_tieptuyen),4)}$}, nếu tiếp tục kẻ tiếp tuyến với hàm $f$ ở điểm này thì ta có điểm rất gần với nghiệm (phương pháp Newton thường hội tụ nhanh trong lân cận của nghiệm)

\cen{\sageplot[scale=.35]{plotall_tieptuyen}}

\newpage\titls{Phương pháp Newton}

%decipoint=5: sẽ có báo lỗi do có sai số kiểu 1x10^(-5), bỏ qua lỗi này được
\begin{sagesilent}
 n=5
 decipoint=4
 xk=range(n+1)
 xk[0]=2
 fxk=range(n)
 f_diff_xk=range(n)
 deltaxk=range(n)
 p_f=find_root(f,0,3) # find nearly exact root
 for i in range(n):
    fxk[i]=f(x=xk[i])
    f_diff_xk[i]=f_diff(x=xk[i])
    xk[i+1]=(f_diff_xk[i]*xk[i]-fxk[i])/f_diff_xk[i]
    deltaxk[i]=abs(xk[i]-p_f)
 xk_disp=[round(xk[i],decipoint) for i in range(n)]
 fxk_disp=[round(fxk[i],decipoint) for i in range(n)]
 deltaxk_disp=[round(deltaxk[i],decipoint) for i in range(n)]
 p_f_disp=round(p_f,decipoint)
 tableN1=r"\begin{tabular}{c|ccc}"
 tableN1+=r"$k$ & $x_k$ & $f(x_k)$ & $|x_k-p|$ \\ \hline"
 for i in range(n):
   tableN1+=latex(i) + r"&" + latex(xk_disp[i]) + r"&" + latex(fxk_disp[i]) + r"&" + latex(deltaxk_disp[i]) + r"\\"
 tableN1+=r"\end{tabular}"
\end{sagesilent}

Ví dụ: tìm nghiệm của $f(x)=\sage{f}$ bằng phương pháp Newton với $x_0=2$ (nghiệm chính xác: $p=\sage{p_f_disp}$)

\cen{\sagestr{tableN1}}

\newpage\titls{Phương pháp Newton}

% Hàm f2 khai báo ở ví dụ về phương pháp lặp đơn

\begin{sagesilent}
 n=9
 decipoint=5
 xk=range(n+1)
 xk[0]=0
 f2xk=range(n)
 f2_diff=diff(f2,x,1)
 f2_diff_xk=range(n)
 deltaxk=range(n)
 p_f2=find_root(f2,0,3) # find nearly exact root
 for i in range(n):
    f2xk[i]=f2(x=xk[i])
    f2_diff_xk[i]=f2_diff(x=xk[i])
    xk[i+1]=(f2_diff_xk[i]*xk[i]-f2xk[i])/f2_diff_xk[i]
    deltaxk[i]=abs(xk[i]-p_f2)
 xk_disp=[round(xk[i],decipoint) for i in range(n)]
 f2xk_disp=[round(f2xk[i],decipoint) for i in range(n)]
 deltaxk_disp=[round(deltaxk[i],decipoint) for i in range(n)]
 p_f2_disp=round(p_f2,decipoint)
 tableN2=r"\begin{tabular}{c|ccc}"
 tableN2+=r"$k$ & $x_k$ & $f(x_k)$ & $|x_k-p|$ \\ \hline"
 for i in range(n):
   tableN2+=latex(i) + r"&" + latex(xk_disp[i]) + r"&" + latex(f2xk_disp[i]) + r"&" + latex(deltaxk_disp[i]) + r"\\"
 tableN2+=r"\end{tabular}"
\end{sagesilent}

Ví dụ: tìm nghiệm của $f(x)=\sage{f2(x)}$ bằng phương pháp Newton với $x_0=0$ (nghiệm chính xác: $p=\sage{p_f2_disp}$)

\cen{\sageplot[scale=.4]{plot(f2,-1,4)}}

\newpage\titls{Phương pháp Newton}

Ví dụ: tìm nghiệm của $f(x)=\sage{f2(x)}$ bằng phương pháp Newton với $x_0=0$ (nghiệm chính xác: $p=\sage{p_f2_disp}$)

\cen{\sagestr{tableN2}}

\newpage\titls{Bài tập}

\itemi{
\item Viết rõ công thức để thực hiện các giải thuật đối với các ví dụ trong bài này (phục vụ tính toán bằng tay)
\item Dùng máy tính cầm tay để thực hiện các giải thuật nhằm giải các ví dụ, so sánh với kết quả trong bài học
\item Làm các bài tập 1-7 trong chương 2, sách "Giáo trình phương pháp tính" (Lê Thái Thanh)
\item Lập trình các thuật toán (thành các file hàm .m trên Octave) theo các phương pháp: 
  \itemi{
  \item PP chia đôi, PP dây cung
  \item PP lặp đơn, PP Newton
  }
}

\label{lastpage}
\end{document}
